\section{Conclusion}

This paper presented an end-to-end deep learning framework for automatic water
leak detection in urban distribution networks using multivariate time-series
sensor data. The proposed approach integrates domain-informed feature
engineering, sliding-window sequence construction, and a hybrid
CNN-BiLSTM-Attention architecture to effectively capture both local and
long-range temporal dependencies.

Experimental results demonstrate that the model achieves strong performance in
detecting leak events, with high AUC-ROC and AUC-PR scores and a favorable
balance between precision and recall. The use of threshold calibration further
enhances the practical applicability of the system by allowing flexible control
over the trade-off between detection sensitivity and false alarm rate. This is
particularly important in real-world monitoring scenarios where operational
constraints must be carefully considered.

The results confirm that combining convolutional, recurrent, and attention-based
mechanisms provides a powerful solution for time-series leak detection. In
addition, the proposed pipeline emphasizes the importance of proper data
preprocessing, leakage-free evaluation, and imbalance-aware metrics, which are
critical for developing reliable and deployable machine learning systems.

Future work will focus on extending the model to more diverse datasets and
network configurations, incorporating baseline comparisons and ablation studies,
and improving interpretability. Furthermore, expanding the framework to support
leak localization and severity estimation represents a promising direction for
enhancing its utility in smart water management systems.

Overall, this work contributes to the development of intelligent and reliable
monitoring solutions for urban water infrastructure, with potential benefits in
reducing water loss, improving operational efficiency, and supporting
sustainable resource management.
